\section{Coverage Roadmap}
\label{sec:roadmap}

The product roadmap is organized around the \textbf{tiered coverage strategy}: achieving
full Mermaid compatibility in priority order, then extending beyond Mermaid.

\subsection{Strategy: Mermaid-Complete First, Net-New After}

The sequencing discipline:
\begin{enumerate}
  \item Cover all 22 Mermaid diagram types (Tiers 0--3).
  \item Establish the aesthetic bar for each type (beat Mermaid's rendering).
  \item Only then invest in net-new diagram types with no Mermaid equivalent.
\end{enumerate}

This ensures the migration story (``drop-in replacement with better output'') is complete
before pursuing differentiation that requires user education.

\subsection{Tier 0: Wire Existing Grammars to Mermaid Syntax}

\textbf{Types:} flowchart, sequence, gantt, timeline, mindmap (5 types).

\textbf{Work required:}
\begin{itemize}
  \item Mermaid syntax parsers for each type (PEG/Nearley grammars)
  \item CST $\to$ Domain IR lowering for each parser
  \item Compatibility test suite (parse Mermaid examples $\to$ verify output)
  \item Frontmatter/directive handling in parser layer
\end{itemize}

\textbf{Kernel work:} None. All four layout engines are built and tested.

\textbf{Exit criteria:} Standard Mermaid examples from \texttt{mermaid.js.org} documentation
parse and render with our compiler, producing visually superior output.

\subsection{Tier 1: The UML/Software Line}

\textbf{Types:} class, state, erDiagram, C4Context, C4Container, C4Component, C4Dynamic
(7 types).

\textbf{Work required:}
\begin{itemize}
  \item New Domain IRs: ClassDiagram IR, StateDiagram IR, ER IR
  \item C4 IR (extends Flow IR with stereotype/boundary semantics)
  \item Layout: class/state use layered layout (reuse Sugiyama kernel); ER uses
        force-directed or layered
  \item Mermaid syntax parsers for all UML types
  \item UML-specific theme extensions (stereotypes, compartments, cardinality labels)
  \item ``Modern UML'' visual standard (beating PlantUML's dated rendering)
\end{itemize}

\textbf{Kernel work:} Minor extensions to node rendering (compartmented boxes, stereotype
headers, cardinality decorators on edges).

\textbf{Exit criteria:} Real-world UML diagrams from popular open-source projects render
beautifully. Side-by-side with PlantUML shows clear aesthetic superiority.

\subsection{Tier 2: The Chart Family}

\textbf{Types:} pie, xychart, quadrant, radar (4 types).

\textbf{Work required:}
\begin{itemize}
  \item New Chart Domain IR and JSON Schema (Section~\ref{sec:chart-family})
  \item Scale computation library (linear, band, radial)
  \item Chart layout engine compiling to Scene IR
  \item Mermaid parsers for all four chart syntaxes
  \item Chart-specific theme extensions (axis styling, gridlines, palettes, legends)
\end{itemize}

\textbf{Kernel work:} The chart layout engine is genuinely new architecture---the only
family that requires a new kernel rather than reusing an existing one.

\textbf{Exit criteria:} All four chart types render from Mermaid syntax with proper axes,
legends, and responsive sizing. Output quality comparable to simple Vega-Lite charts.

\subsection{Tier 3: Remaining Types}

\textbf{Types:} sankey, requirement, gitGraph, block, architecture, treemap, kanban,
journey, packet (9 types).

\textbf{Work required per type:}
\begin{itemize}
  \item Mermaid syntax parser
  \item Domain IR definition (most are minor variants of existing IRs)
  \item Layout integration (most reuse node-link, tree, or timeline kernels)
  \item Type-specific rendering (sankey flow widths, git branch colors, etc.)
\end{itemize}

\textbf{Kernel work:} Minimal. Sankey requires weighted-edge rendering; treemap requires
squarified rectangle packing; others reuse existing layout families.

\textbf{Exit criteria:} All 22 Mermaid diagram types render correctly.

\subsection{MVP Definition (Reframed)}

The MVP is defined as:

\begin{description}
  \item[Tier 0 complete] All five existing grammars accessible via Mermaid syntax parsers.
  \item[Aesthetic bar met] Output demonstrably superior to Mermaid for all Tier-0 types.
  \item[Agent IR path working] Structured IR input for all Tier-0 types with published
        JSON Schemas.
  \item[CLI + library] \texttt{timeline render input.mmd -o output.svg} works.
  \item[Composition] Multi-diagram posters from superset DSL.
\end{description}

The MVP answers: \emph{Can users drop in Mermaid files and get better output, with zero
migration cost, plus the agent IR and composition capabilities Mermaid lacks?}

\subsection{Post-MVP: Net-New Diagram Types}

After Mermaid-complete coverage, new diagram types with no Mermaid equivalent:
\begin{itemize}
  \item \textbf{Architecture Decision Records} (ADR visualization)
  \item \textbf{Dependency graphs} (package/module dependency visualization)
  \item \textbf{Infrastructure diagrams} (cloud architecture beyond C4)
  \item \textbf{Data pipeline diagrams} (ETL/streaming topology)
\end{itemize}

These are post-Tier-3 and driven by user demand, not pre-planned.

\subsection{Implementation Status Summary}

\begin{table}[h]
\centering\small
\caption{Current implementation status vs.\ roadmap tiers}
\begin{tabular}{llll}
\toprule
\textbf{Component} & \textbf{Tier} & \textbf{Status} & \textbf{Tests} \\
\midrule
Timeline grammar + themes   & 0 & Built & 551 \\
Flow grammar + themes       & 0 & Built & 33 \\
Sequence grammar + themes   & 0 & Built & 37 \\
Tree grammar + themes       & 0 & Built & 26 \\
Composition layer           & --- & Built & 25 \\
Animation (SMIL)            & --- & Built & incl.\ in flow \\
Schema validation           & --- & Built & 13 \\
Mermaid DSL parsers         & 0 & \textbf{Planned} & --- \\
UML/software line           & 1 & Planned & --- \\
Chart family                & 2 & Planned & --- \\
Remaining types             & 3 & Planned & --- \\
\bottomrule
\end{tabular}
\label{tab:implementation-status}
\end{table}

\textbf{Total tests passing:} 795 (all grammars + composition + schema + CLI).
All existing goldens byte-identical across runs.
